\documentclass[twoside,openany]{book}
\usepackage[czech]{babel}          % Language
\usepackage[utf8]{inputenc}        % Encoding of characters in this .tex file
\usepackage{cmap}                  % Make PDF file searchable and copyable (ASCII characters)
\usepackage{lmodern}               % Make PDF file searchable and copyable (Accented characters)
\usepackage[T1]{fontenc}           % Hyphenate accented words
\usepackage[a6paper, top=0.75cm, bottom=1.5cm, inner=0.75cm, outer=0.75cm]{geometry}  % Paper size and margins
\usepackage[fontsize=9]{scrextend} % Font size
\usepackage{marginnote}            % Enable margin notes
\usepackage{rotating}              % Rotated environments (for margin headers)
\usepackage[protrusion]{microtype} % Better typeset results
\usepackage{enumitem}              % Nicer enumeration lists
% Hack follows: my version of csquotes does not know Czech (but Danish style is the same)
\usepackage[style=danish]{csquotes} % Add command to enquote strings
\usepackage{ragged2e}              % Add \justify command
\usepackage{titlesec}              % Customize chapters
\usepackage[usenames,dvipsnames,svgnames,table]{xcolor} % Custom colors
\usepackage{tikzsymbols}           % Smileys
\usepackage{chronology}            % Timeline
\usepackage{tikz}                  % Number lines
\usepackage{url}                   % Enable URL typesetting

% ======= CUSTOM DOCUMENT SETTINGS ========

% Specify relevant Tikz libraries
\usetikzlibrary{arrows}

% Change \chapter and \section command display and spacings
% Source: https://tex.stackexchange.com/questions/63390/how-to-decrease-spacing-before-chapter-title
\titleformat{\chapter}[display]{\normalfont\Large\bfseries}{Týden \thechapter}{0em}{\Large}
\titlespacing*{\chapter}{0em}{-2em}{0em}
\titleformat{\section}[display]{\normalfont\bfseries}{\sectionname}{}{}
\titlespacing*{\section}{0em}{0.5em}{0em}

% \punct command allows to shift the footnotemark above a punctuation
\newlength{\spc} % declare a variable to save spacing value
\newcommand{\punct}[1]{%
  \settowidth{\spc}{#1}% set value of \spc variable to the width of #1 argument
  \addtolength{\spc}{-1.8\spc}% subtract from \spc about two (1.8) of its values making its magnitude negative
  #1% print the optional argument
  \hspace*{\spc}% print an additional negative spacing stored in \spc after #1
}

% Remove headers and position page number in the bottom mid
\pagestyle{plain}

% Smaller spaces between enumeration lists
\setitemize{itemsep=0em}

% Disable automatic indentation, enable paragraph skip
\setlength{\parindent}{0em}
\setlength{\parskip}{0.5em}


% Rule length including the margin
\newlength{\rulelength}
\setlength{\rulelength}{\dimexpr(\textwidth+0.5cm)\relax}

% Smaller note (italics, gray)
\newcommand{\note}[1]{\textcolor{darkgray}{\small\itshape #1}}
\newcommand{\noteblack}[1]{{\small\itshape #1}}

% accurate absolute positioning
\newcommand{\putat}[3]{\begin{picture}(0,0)(0,0)\put(#1,#2){#3}\end{picture}}

% ======= METADATA =======

\title{Reflektivní deník učitele}
\author{Kolektiv vyučujících Praktika vedení cvičení}
\date{\today}

% ======= DOCUMENT START =======

\begin{document}

% Front cover page
\thispagestyle{empty}
\putat{-21}{-390}{\includegraphics[width=\paperwidth]{../img_cover/cover-front}}
\newpage
\thispagestyle{empty}

Tato verze Reflektivního deníku učitele je upravená pro potřebu učitelů a vedoucích kroužků na Cyrilometodějském gymnáziu a střední odborné škole pedagogické Brno. Původní verze je určena cvičícím na Fakultě informatiky Masarykovy univerzity.

Úpravy činí: dva týdny navíc (do celkových \textbf{šestnácti}), logo školy na přebalu a tento úvod.

Terminologie ve zbytku deníku zůstává nezměněná, je tedy potřeba text převádět do vlastního kontextu, předně:

\begin{itemize}	
\item    cvičící -> učitel / vedoucí kroužku
\item    semester -> pololetí
\item    zkouška -> pololetní písemka / závěrečný test
\item	...
\end{itemize}
	
Deník je doplňkem semináře pro učitele na FI MUNI jménem \textit{Teaching lab}, ve kterém se zavádí různé pojmy a koncepty, jejichž výčet se nachází na konci tohoto deníku ("Užitečné výukové nástroje"). Jeden koncept však stojí za dovysvětlení již teď, neboť se nachází hned v reflexi prvního týdne:

\textbf{precedens} (2. p., č. j. bez precedentu, 1. p., mn. č. precedenty) je něco, co se v hodině stane poprvé a studenti pak mají tendenci to tak opakovat (např. že se mluví bez přihlášení nebo že se toleruje chození pozdě). Jinými slovy: "Co se stane jednou, může se stát dvakrát. Co se stane dvakrát, stane se potřetí."

% Foretitle
\begin{titlepage}
	\centering
	\vspace*{1cm}
	\includegraphics[width=0.5\textwidth]{../img/diary}\par
	\vspace{0.5cm}
	{\huge\bfseries Reflektivní\\ deník učitele\par}
	\vfill
	\raggedright
	\Large
	Majitel:\par
	\vspace{0.2cm}
	Předmět:\par
\end{titlepage}

% ===== DIARY CONTENTS =====

% ===== INTRODUCTION =====

\chapter*{Jak se z člověka stane dobrý učitel?}

\vspace*{1em}
Položme si namísto statické otázky dynamickou:

\vspace*{1em}
\textit{\large \enquote{Co mohu dělat pro to, abych se jako učitel zlepšoval?}}

\vspace*{1em}
Ukazuje se, že pro zlepšení jsou nutné (a nejspíš i dostačující) tři věci\punct{.}\footnotemark
\footnotetext{Jsou i další věci, které mohou pomoci (například expertní zpětná vazba), ale ty nejsou dobře škálovatelné.}
\begin{enumerate}
\item \textbf{Pravidelně učit.}\\Potřebuješ mít vlastní zkušenost, nejlépe často a pravidelně opakovanou.
\item \textbf{Reflektovat svoje učení.}\\Všímej si a přemýšlej nad tím, co fungovalo dobře a co příště zkusíš jinak.
\item \textbf{Pozorovat, jak učí ostatní.}\\Přemýšlej o tom, co dělají ostatní učitelé, co jim funguje, co ne a co z toho můžeš použít.
\end{enumerate}

Tento deníček ti pomůže reflektovat způsob, jakým učíš. Nalezneš v~něm informace o tom, čeho si všímat, jaké otázky si klást a nad čím přemýšlet. Není to však učebnice dobrého učitele -- je to pouze průvodce na cestě, po které už jdeš.

Vhodným doplňkem k vlastní reflexi výuky jsou diskuze s ostatními učiteli o tom, jak učí. Pokud se chceš setkávat s dalšími cvičícími pravidelně a navíc se dozvědět více o tom, co funguje a co ne, zapiš si předmět \textit{Teaching Lab} na Fakultě informatiky Masarykovy univerzity.

\newpage
\section*{Jak používat reflektivní deník?}

Především je důležité používat ho pravidelně.

Zapiš si svoje myšlenky při přípravě hodiny i po odvedené výuce. Máš nachystaných čtrnáct dvoustránek (jednu na každý týden semestru), přičemž každá nabízí pár otázek. Deníček je malý -- to mimo jiné i proto, aby ti psaní nezabralo moc času a tvé poznámky byly krátké.

Za stránkami pro jednotlivé týdny najdeš ještě rubriku učitelských schopností, seznam indikátorů, které je možné na hodině sledovat, seznam některých užitečných nástrojů k výuce a na závěr prostor na své vlastní poznámky.

Někde záměrně používáme anglické termíny (jako například \textit{tracking}), protože vhodný český překlad buď neexistuje, nebo má konotace, které nechceme.

\section*{Proč používat reflektivní deník?}

Pravidelná práce s deníkem plní vícero cílů:
\begin{enumerate}[topsep=0pt]
\item Pomáhá ti zamyslet se nad svou hodinou (reflexe).
\item Nabízí prostor pro snadné psaní si poznámek pro příště.
\item Pomáhá ti uvědomit si, co všechno patří k tvojí roli učitele.
\item Umožňuje ti sledovat svůj vlastní pokrok.
\end{enumerate}

\section*{Máš nápady na zlepšení?}

Budeme velice rádi, když nám sdělíš svoje návrhy na zlepšení nebo postřehy z používaní deníku. Napiš nám email na adresu \textit{teachinglab@fi.muni.cz}.

% ===== REFLECTION ON INDIVIDUAL WEEKS =====

% Margin note
\newcommand{\marginsection}[3]{\marginpar{\raisebox{#2\height}{
\begin{turn}{#1}
\bfseries \color{gray} \large
#3\end{turn}}}}

% Timeline for planning the seminar structure
\newcommand{\timeline}{
\vspace*{2.1cm}
\color{gray}
\begin{chronology}[15]{0}{120}{\textwidth}
\end{chronology}
\vspace*{-0.2cm}
\color{black}
}

% Smileys for evaluating teacher's own feelings
\newcommand{\smileys}{
Jak jsem spokojen(a) s hodinou?
\raisebox{-0.3em}{\,\,\Annoey[2]\,\,\Sadey[2]\,\,\Changey[2]{0}\,\,\Changey[2]{1}\,\,\Changey[2]{2}}
}

% Summarizing 2 good and 2 bad points of the week
\newcommand{\goodbadpoints}{
Co se mi povedlo?
\begin{enumerate}[topsep=0pt,itemsep=1pt]
\item
\item
\end{enumerate}

Co mohlo být lépe?
\begin{enumerate}[topsep=0pt,itemsep=1pt]
\item
\item
\end{enumerate}
}

% Common week
\newcommand{\commonweek}{
\chapter{}
\vspace*{-2em}

\marginsection{90}{-1}{před výukou}

Příprava:\hspace{1cm} hodin(y)

Jaké jsou cíle této hodiny? Proč chci studenty učit tohle?
\vspace*{1cm}

Jak poznám, že se cílů podařilo dosáhnout?

\hspace*{-1cm}
\rule{\rulelength}{0.4pt}

\marginsection{90}{-3}{po výuce}

\smileys

\goodbadpoints

Vlastní otázky:

\newpage
\marginsection{-90}{1}{volné poznámky}
}

\newgeometry{top=0.75cm, bottom=1.5cm, inner=0.75cm, outer=1.2cm,
			heightrounded, marginparwidth=0.5cm, marginparsep=0.3cm}

% ======= WEEK 1 =======

\chapter{}
\vspace*{-2em}

\marginsection{90}{-1.3}{vyplnit před výukou}

Příprava:\hspace{1cm} hodin(y)\\
\note{Příprava materiálů, úkolů, pokynů, \ldots}

Jaká je struktura mojí hodiny?\\
\note{Rozvrhni si 2--6 bloků na časovou osu níže.}

\timeline

Jaké precedenty chci ve výuce nastavit? \\
\note{Oslovování, (ne)formálnost, kladení otázek, začínání včas, \ldots}
\vspace*{1cm}

\hspace*{-1cm}

\rule{\rulelength}{0.4pt}

\marginsection{90}{-1.3}{vyplnit po výuce}

\smileys


\note{Jaký mám pocit já? Co vidím na studentech?}

\goodbadpoints

\newpage
\marginsection{-90}{1}{volné poznámky}
\vspace*{-2em}
\note{Jaké (ne)chtěné precedenty se objevily na hodině?\\
Jak na další hodině udržím ty chtěné a napravím ty nechtěné?\\
Jaká byla atmosféra během vyučování?\\
Rozumí studenti organizaci a struktuře kurzu?\\
Ví studenti, co se od nich očekává?\\
(Další tipy na otázky najdeš na straně \pageref{indicators}.)}

% ======= WEEK 2 =======

\chapter{}
\vspace*{-2em}

\marginsection{90}{-1.2}{vyplnit před výukou}

Příprava:\hspace{1cm} hodin(y)\\
\note{Příprava materiálů, úkolů, pokynů, ...}

Jaká je struktura mojí hodiny?\\
\note{Rozvrhni si 2--6 bloků na časovou osu níže.}

\timeline

Jaké precedenty chci ve výuce udržet? \\
\note{Oslovování, (ne)formálnost, kladení otázek, začátek včas, ...}
\vspace*{1cm}

\hspace*{-1cm}

\rule{\rulelength}{0.4pt}

\marginsection{90}{-1.3}{vyplnit po výuce}

\smileys
\note{Jaký mám pocit já? Co vidím na studentech?}

\goodbadpoints

\newpage
\marginsection{-90}{1}{volné poznámky}
\vspace*{-2em}
\note{Na jaký aspekt výuky se zaměřím při nejbližší hospitaci?\\
Jaký aspekt výuky má sledovat kolega, který navštíví mou hodinu?\\
(Další tipy na otázky najdeš na straně \pageref{indicators}.)}

% ======= WEEK 3 =======

\chapter{}
\vspace*{-2em}

\marginsection{90}{-1.1}{vyplnit před výukou}

Příprava:\hspace{1cm} hodin(y)

Jaká je struktura mojí hodiny?

\timeline

Jaké otázky na hodině položím skupině?
\vspace*{0.8cm}


\hspace*{-1cm}
\rule{\rulelength}{0.4pt}

\marginsection{90}{-1.5}{vyplnit po výuce}

\smileys

\goodbadpoints

Na které otázky skupina reagovala? Na které ne?

\newpage
\marginsection{-90}{1}{volné poznámky}
\vspace*{-2em}
\note{Kolik otázek do publika jsem položil(a)?\\
Kolik otázek položili studenti mně?\\
Je to málo, tak akorát, nebo příliš?\\
(Další tipy na otázky najdeš na straně \pageref{indicators}.)}

% ======= WEEK 4 =======

\chapter{}
\vspace*{-2em}

\marginsection{90}{-1}{vyplnit před výukou}

Příprava:\hspace{1cm} hodin(y)

Jaká je struktura mojí hodiny?

\timeline

\hspace*{-1cm}
\rule{\rulelength}{0.4pt}

\marginsection{90}{-1.6}{vyplnit po výuce}

\smileys

\goodbadpoints

Vlastní otázka:\\
\note{Sem napiš otázku, na kterou si chceš po výuce odpovědět.}

\newpage
\marginsection{-90}{1}{volné poznámky}
\vspace*{-2em}
\note{Jaké znalosti a dovednosti chci naučit? Popiš je ve formátu rubriky, tedy pojmenuj dovednost a popiš škálu (viz str.\ \pageref{rubric}).\\
Jaký typ učitele chci být? Co bych měl(a) pro to dělat?\\
(Další tipy na otázky najdeš na straně \pageref{indicators}.)}

% ======= WEEK 5 =======

\chapter{}
\vspace*{-2em}

\marginsection{90}{-1.2}{před výukou}
Příprava:\hspace{1cm} hodin(y)

Jaké jsou cíle této hodiny? Proč chci studenty učit tohle?\\
\note{Při formulování cílů ti pomůžou slovesa na straně \pageref{bloom}.}
\vspace*{1cm}

Jak poznám, že se cílů podařilo dosáhnout?\\
\note{Existuje způsob, kterým je možné zjistit, zda jsi cíle dosáhl(a)?}

\hspace*{-1cm}
\rule{\rulelength}{0.4pt}

\marginsection{90}{-2.5}{po výuce}

\smileys

\goodbadpoints

Vlastní otázka:\\
\note{Sem napiš otázku, na kterou si chceš po výuce odpovědět.}

\newpage
\marginsection{-90}{1}{volné poznámky}
\vspace*{-2em}
\note{Podařilo se mi naplnit cíle hodiny?\\
Co je třeba do budoucna na téhle hodině změnit?\\
Na co bych neměl zapomenout, až budu plánovat další hodinu?\\
(Další tipy na otázky najdeš na straně \pageref{indicators}.)}

% ======= WEEK 6 =======

\chapter{}
\vspace*{-2em}

\marginsection{90}{-1.2}{před výukou}

Příprava:\hspace{1cm} hodin(y)

Jaké jsou cíle této hodiny? Proč chci studenty učit tohle?\\
\note{Při formulování cílů ti pomůžou slovesa na straně \pageref{bloom}.}
\vspace*{1cm}

Jak poznám, že se cílů podařilo dosáhnout?\\
\note{Existuje způsob, kterým je možné zjistit, zda jsi cíle dosáhl(a)?}

\hspace*{-1cm}
\rule{\rulelength}{0.4pt}

\marginsection{90}{-2.5}{po výuce}

\smileys

\goodbadpoints

Vlastní otázka:\\
\note{Sem napiš otázku, na kterou si chceš po výuce odpovědět.}

\newpage
\marginsection{-90}{1}{volné poznámky}
\vspace*{-2em}
\note{Podařilo se mi naplnit cíl hodiny?\\
Co je třeba do budoucna na téhle hodině změnit?\\
Co je problematické na úkolech, které jsem zadal?\\
(Další tipy na otázky najdeš na straně \pageref{indicators}.)}

% ======= WEEK 7 =======
\commonweek

% ======= WEEK 8 =======
\commonweek

% ======= WEEK 9 =======
\commonweek

% ======= WEEK 10 =======
\commonweek

% ======= WEEK 11 =======
\commonweek

% ======= WEEK 12 =======
\commonweek

% ======= WEEK 13 =======
\commonweek

% ======= WEEK 14 =======
\commonweek

% ======= WEEK 15 =======
\commonweek

% ======= WEEK 16 =======
\commonweek

% ======= WEEK 17 =======

\chapter{}
\vspace*{-2em}

\marginsection{90}{-1}{před výukou}

Příprava:\hspace{1cm} hodin(y)

Jaké jsou cíle této hodiny? Proč chci studenty učit tohle?
\vspace*{1cm}

Jak poznám, že se cílů podařilo dosáhnout?

\hspace*{-1cm}
\rule{\rulelength}{0.4pt}

\marginsection{90}{-3}{po výuce}

\smileys

\goodbadpoints

Blíží se konec semestru -- dostávají studenti znalosti a dovednosti, které potřebují ke zkoušce a pro praxi?\\
\note{Jestli ne, jaké změny můžu v následujících týdnech provést?}

\newpage
\marginsection{-90}{1}{volné poznámky}

% ======= WEEK 18 =======

\chapter{}
\vspace*{-2em}

\marginsection{90}{-1}{před výukou}

Příprava:\hspace{1cm} hodin(y)

Jaké jsou cíle této hodiny? Proč chci studenty učit tohle?
\vspace*{1cm}

Jak poznám, že se cílů podařilo dosáhnout?

\hspace*{-1cm}
\rule{\rulelength}{0.4pt}

\marginsection{90}{-3}{po výuce}

\smileys

\goodbadpoints

Jaký je širší kontext toho, co učím?\\
\note{Kde v mojí výuce se dotýkám jiných kurzů nebo jiných oblastí vědění?}

\newpage
\marginsection{-90}{1}{volné poznámky}
\vspace*{-2em}
\note{Jak získám od studentů zpětnou vazbu na ty aspekty mé vyúky, které mě zajímají?\\
Jak to dělají ostatní vyučující v mém předmětu?\\
A jak dělají ostatní vyučující v jiných předmětech?}

% ======= WEEK 19 =======

\chapter{}
\vspace*{-2em}

\marginsection{90}{-1}{před výukou}

Příprava:\hspace{1cm} hodin(y)

Jaké jsou cíle této hodiny? Proč chci studenty učit tohle?
\vspace*{1cm}

Jak poznám, že se cílů podařilo dosáhnout?

\hspace*{-1cm}
\rule{\rulelength}{0.4pt}

\marginsection{90}{-3}{po výuce}

\smileys

\goodbadpoints

Jak jsem pokročil(a) za tento semestr výuky?\\
\note{Co jsem si uvědomil(a)? Co jsem zlepšil(a)?}

\newpage
\marginsection{-90}{1}{volné poznámky}
\vspace*{-2em}
\note{Za jaké znalosti a dovednosti jsou studenti hodnoceni u zkoušky?\\
Odpovídají metody hodnocení tomu, co učím?\\
Pokud mají studenti zájem se dál samostatně rozvíjet v dané oblasti, jsou na to po tomto kurzu připraveni?}

% ======= WEEK 20 =======

\chapter{}
\vspace*{-2em}

\marginsection{90}{-1}{před výukou}

Příprava:\hspace{1cm} hodin(y)

Jaké jsou cíle této hodiny? Proč chci studenty učit tohle?
\vspace*{1cm}

Jak poznám, že se cílů podařilo dosáhnout?

\hspace*{-1cm}
\rule{\rulelength}{0.4pt}

\marginsection{90}{-3}{po výuce}

\smileys

\goodbadpoints

Co by se mělo další semestr změnit?\\
\note{Co pro mě nebo studenty nefungovalo?}\\
Co by mělo další semestr zůstat stejné?

\newpage
\marginsection{-90}{1}{volné poznámky}
\vspace*{-2em}
\note{Co fungovalo pro mé kolegy?\\
Můžeme si sdílet know-how, jak učit lépe?}

% ===== RUBRIC OF TEACHER'S JOB ASPECTS =====

% A single rubric criterion
\newcommand{\rubriccriterion}[4]{
\stepcounter{rubricquestion}
\section*{\therubricquestion: #1}

\smallskip
\note{Neznalý:} #2

\note{Začátečník:} #3

\note{Guru:} #4

\medskip
\begin{tikzpicture}
\draw (0,0) -- (8,0);
\foreach \i in {0,1,...,8} % numbers on line
{
\fill[black] (\i,0) circle (1.5 mm);
\fill[white] (\i,0) circle (1.4 mm);
}
\node at (0.15, -0.5) {neznalý};
\node at (3, -0.5)    {začátečník};
\node at (8, -0.5)    {guru};
\end{tikzpicture}
}

\restoregeometry
\chapter*{Učitelská rubrika}
\label{rubric}

Následující strany prezentují rubriku činností učitele. Běžně si učitelé-začátečníci (a někdy ani učitelé-veteráni) neuvědomují všechny úrovně a dimenze učitelského umění. To se často stává proto, že se s tím nikdy nesetkali u ostatních. Je možné, že po pár měsících či letech učení přijde moment, kdy si uvědomíte, že máte mnohem větší prostor pro zlepšování, než jste si prvně mysleli\punct{.}\footnote{Více informací najdete například v knize \emph{How Learning Works: Seven Research-Based Principles for Smart Teaching}.}

\section*{Co je to rubrika?}

Rubrika je nástroj pro \emph{self-assessment} (uvědomění a popsání svých vlastních znalostí a dovedností). Pomáhá zřetelně si uvědomit a pojmenovat, jaké části daný koncept má (např.\ umět programovat znamená být schopen navrhnout algoritmus, efektivně pracovat s~pamětí, znát syntaxi, \dots). Zároveň umožňuje vyhodnotit vlastní úroveň v jednotlivých částech (např.\ navrhnout algoritmus umím dobře, zapsat to kódem mám trochu problém a práci s pamětí neoptimalizuji vůbec).

\section*{Jak učitelskou rubriku používat?}

Na začátku semestru si rubriku vyplň (označ, kde na dané škále jsou podle tebe tvoje schopnosti). Vyber si 1--3 oblasti, na které se chceš zaměřit v tomto semestru. Chceš-li být hodně důsledný(á), můžeš vymyslet i konkrétní kroky, které uděláš, a indikátory svého postupu. Na konci semestru si můžeš rubriku vyplnit znovu a vyhodnotit, v čem jsi se posunul(a).

Rubriku můžeš čist i jako manifest -- popisy kategorie \enquote{guru} vystihují naši představu o tom, co zvládají skvělí učitelé a k čemu bychom rádi vedli začínající učitele.

\newcounter{rubricquestion}

\newpage
\rubriccriterion{Cíle hodiny, vědomá pozornost}
{Ve výuce nemám jasně stanovené cíle. Když učím, nevnímám, co se děje a kam směřujeme. Často se cítím ztracen(a).}
{Někdy se mi daří uvědomovat si, jaký mám cíl, co právě dělám a jaký to bude mít efekt. Většinou ale nemám při výuce vědomou pozornost.}
{Během výuky jsem si téměř vždy vědom(a) toho, jaký cíl sleduji, co se se mnou i se skupinou děje, co dělám a jaký to bude mít efekt. Vím, jak jsme se dostali tam, kde právě jsme.}

\rule{\textwidth}{0.4pt}
\rubriccriterion{Interakce se skupinou, kladení otázek skupině}
{Se skupinou neinteraguji. Neptám se, protože bych asi nedostal(a) odpověď. Nevím, jak studenty zapojit.}
{Vím, že lze interagovat se skupinou a znám nástroje (chápu, jak by to šlo). Nedokážu je ale dobře používat. Někdy  se skupiny ptám, ale nedostávám odpověď.}
{Se skupinou často a efektivně interaguji. Ptám se způsobem, který studenty aktivizuje a zapojuje. Situace, kdy nedostávám od skupiny odpověď dokážu obratně řešit (např.\ přeformulováním otázky).}

\newpage
\rubriccriterion{Strukturování výuky}
{O strukturování výuky vůbec nepřemýšlím.}
{Chápu smysl přehledného strukturování své hodiny a snažím se o to. Často se ale zamotám, ztratím nit, nebo řeším více věcí naráz a studenti se pak ztrácí nebo odpojují.}
{Moje hodiny mají jasnou strukturu. Studenti ví, co se právě děje, co bude následovat a chápou návaznosti. Mezi jednotlivými bloky vědomě dělám zřetelné přechody.}

\rule{\textwidth}{0.4pt}
\rubriccriterion{Vlastní pocity a spokojenost s výukou}
{Jsem nespokojen(a) se svými hodinami anebo své pocity a vlasní spokojenost s výukou nijak nereflektuji. Netěším se do výuky.}
{Ve výuce si často moc nevěřím. Někdy mi vyjdou části hodiny, ale často se cítím napjatě nebo mám strach, že něco pokazím.}
{Ve výuce se cítím uvolněně a sebevědomě. Baví mě to a mám svůj styl.}

\newpage
\rubriccriterion{Formativní zpětná vazba}
{O dávání osobní zpětné vazby studentům nijak nepřemýšlím.}
{Snažím se studentům zpětnou vazbu dávat, aby mohli růst. Myslím si však, že jí není dost, nebo to nedělám efektivně. Studenti mou zpětnou vazbu někdy nevnímají jako podporu a projev respektu.}
{Se studenty ve výuce interaguji tak, že dostávají průběžně formativní zpětnou vazbu. Chápou tedy, co jim jde, kde dělají chyby a jak se mohou zlepšovat. Zároveň cítí, že je respektuji a zpětné vazby se nebojí.}

\rule{\textwidth}{0.4pt}
\rubriccriterion{Jasné zadávání pokynů a úloh}
{O tom, jak zadávat pokyny nebo úlohy, nijak zvlášť nepřemýšlím.}
{Stává se mi, že zadám pokyn nebo úlohu a studenti neví, co dělat, jak začít nebo k čemu mají dojít (co má být výsledkem).}
{Když zadávám pokyn nebo úlohu, studenti mají jasno v~tom, co dělat, jak začít nebo k čemu mají dojít.}

\newpage
\rubriccriterion{Variabilita a inovace ve výuce}
{Učím tak, jak mi řekli nebo kopíruji výuku, kterou jsem zažil(a). Nepřemýšlím o~jiných variantách.}
{Jsem si vědom(a) toho, že existuje mnoho typů aktivit, které lze při výuce použít. Neznám jich ale dostatek, nedokážu je efektivně zadávat, nebo nemám jasno v tom, proč je používat.}
{Znám mnoho různých aktivit a svou výuku skládám tak, aby byla dostatečně pestrá. Vybrané aktivity efektivně procvičují to, co mám záměr procvičovat. Zároveň studenty efektivně zapojují a zvyšují jejich motivaci se učit.}

\rule{\textwidth}{0.4pt}
\rubriccriterion{Širší kontext mé výuky}
{O širším kontextu výuky a mého kurzu nepřemýšlím.}
{Je pro mě těžké jasně pojmenovat znalosti a dovednosti, které u~studentů rozvíjím. Nevím, v jakém kontextu je studenti využijí. Nevidím propojení s dalšími kurzy.}
{Mám jasnou představu o tom, k čemu studenty vedu (jaké dovednosti rozvíjím, jaké znalosti chci předat). Vím, proč tyto dovednosti rozvíjím a v jakém kontextu je studenti v budoucnu použijí. Vidím širší kontext.}
\vspace*{-1em}

\newpage
\rubriccriterion{Jasné vysvětlování}
{Svoje způsoby vysvětlování nijak nereflektuji.}
{Když vysvětluji, běžně se mi stává, že si nejsem jistý, zda vysvětluji dobře a zda to studentům pomáhá při pochopení.}
{Když vysvětluji teorii, demonstruji řešení a efektivně opravuji chyby. Dokážu se dobře vžít do toho, jak to vidí student. Moje vysvělování efektivně pomáhá studentům pochopit problém. Nestává se mi, že bych vysvětloval(a) něco, na co se neptali.}

\rule{\textwidth}{0.4pt}
\rubriccriterion{Nastavení prostředí, systémy ve výuce}
{O atmosféře ve výuce nepřemýšlím. Systémy ve své výuce nevnímám.}
{Přemýšlím o nastavení pravidel i atmosféry. Systémy ve výuce (např.\ bodování) přebírám od ostatních. Nemám ale jasno v jejich efektech, nebo neznám způsoby, jak je lze upravit.}
{Umím ve výuce vytvořit prostředí, které podporuje efektivní učení. U systémů, které používám (např.\ bodování, bonbóny, zahajovací rituály) chápu efekt. Systémy nepřebírám slepě, chápu jejich účel a uzpůsobuji podle svých potřeb.}

\newpage
\rubriccriterion{Flexibilita, přizpůsobování výuky na místě}
{Na situace vzniklé v průběhu výuky nijak vědomě nereaguji.}
{Uvědomuji si chvíle, kdy by mohlo být zajímavé nebo užitečné dělat něco jiného, než jsem měl(a) v plánu. Většinou ale nedokážu v daný moment vhodně zareagovat.}
{Dokážu svoji výuku průběžně přizpůsobovat tomu, co se právě děje ve skupině a co studenti potřebují. Mám k tomu dostatek nástrojů a dokážu je efektivně použít.}

\rule{\textwidth}{0.4pt}
\rubriccriterion{Individuální interakce se studenty}
{O konzultacích nebo jiných indidivuálních interakcích nijak vědomě nepřemýšlím.}
{Stává se mi, že si často nevím rady při individuální interakci se studentem (např.\ zkoušení u~tabule, konzultace). Interakce neprobíhá efektivně nebo se student cítí zastrašen.}
{Při interakcích s jednotlivci (např.\ zkoušení, konzultace) efektivně využívám čas. Studenti se mnou konzultují rádi a je to pro ně užitečné. Cítí ze mě respekt a podporu.}

\newpage
\rubriccriterion{Skupinová práce, management skupiny}
{Přijde mi, že dělit skupinu je zbytečnost. O~práci v~menších skupinkách nepřemýšlím.}
{Uvědomuji si možnosti práce ve dvojicích či menších skupinách. Tuším, že bych toho mohl(a) lépe využívat, hledám jak.}
{Mám jasno v tom, kdy chci pracovat s celou skupinou, kdy s jednotlivci a kdy se skupinkami. Rozdělení do menších skupin efektivně používám, když je to potřeba. Ve vhodných případech zadávám interakce mezi skupinami.}

\rule{\textwidth}{0.4pt}
\rubriccriterion{Čtení atmosféry ve skupině}
{Skupinu ve své výuce nesleduji, pozornost věnuji pouze obsahu.}
{Uvědomuji si, že mi skupina vysílá signály a že by bylo dobré jim rozumět a využít je pro efektivní vedení výuky. Ve výuce to ale dokážu jen výjimečně.}
{Dokážu dobře odhadnout naladění skupiny. Mám jasno v tom, co se ve skupině studentů děje (např. únava, nadšení, obava).}

% ===== OUTRO (INDICATORS, BLOOM, TOOLS) =====

\chapter*{Indikátory}
\label{indicators}
\vspace{-0.5em}
\note{Ukazatel stavu; aspekt výuky, který je možné přímo sledovat.}

\section*{Kvantitativní indikátory}

\textsc{Práce s časem}\\
\noteblack{O kolik minut jsem prodlužoval(a)? Kolik času zůstalo navíc?\\
Jak dlouho trval můj monolog?\\
Kolik času byli studenti aktivně zapojeni?\\
Kolik času jsem v průměru věnoval(a) jednomu studentovi?}

\textsc{Interakce se studenty}\\
\noteblack{Kolik otázek jsem položil(a)? Kolik z nich bylo uzavřených/otevřených?\\
Kolik otázek jsem dostal(a) od studentů?\\
Kolik studentů jsem pochválil(a)/ocenil(a) (za jakkoliv malý úspěch)?}

\textsc{Řešení úloh}\\
\noteblack{Kolik úloh vyřešili studenti na hodině? Kolik úkolů dostali na doma?\\
Kolik studentů bylo ztraceno (těžký úkol)? Kolik se naopak nudilo?}

\textsc{Účast}\\
\noteblack{Kolik studentů přišlo na výuku? Kolik z nich přišlo pozdě?\\
Kolik studentů odešlo během výuky?}

\section*{Kvalitativní indikátory}

\textsc{Pocity a subjektivní dojmy}\\
\noteblack{Jaké pocity jsem měl(a) během výuky?\\
Jak se chovali studenti? Byli zapojení/aktivní?}

\textsc{Struktura výuky}\\
\noteblack{Věděli studenti, co přesně měli dělat, jak, a proč?\\
Dodržel(a) jsem plán hodiny? Pokud ne, proč?}

\textsc{Obsah výuky}\\
\noteblack{Je moje výuka variabilní (různé typy úloh, pomůcek, \dots)?\\
Odpovídá obsah mé výuky tomu, co chci naučit?}

\newpage

\chapter*{Revidovaná Bloomova taxonomie}
\label{bloom}
\vspace{-0.5em}
\note{Hierarchie kognitivních vzdělávacích cílů, B.\ Bloom, 1956\\
Revize D.\ Krathwohl, 2002 (přidána i druhá dimenze)}
\vspace{-0.3em}

\begin{enumerate}[leftmargin=*]
\item \textbf{Vzpomenout si (\textit{remember})}\\
\note{termíny a fakta, jejich klasifikace a kategorizace}\\
{\small definovat, vyjmenovat, zopakovat, popsat, rozpoznat, přiřadit pojem ke konceptu}

\item \textbf{Pochopit (\textit{understand})}\\
\note{překlad do jiné formy, jednoduchá interpretace, extrapolace}\\
{\small interpretovat, reprezentovat, vysvětlit, parafrázovat, sumarizovat, uvést příklad, ilustrovat na příkladu}

\item \textbf{Použít (\textit{apply})}\\
\note{užití postupu ve správné situaci, použití abstrakcí a zobecnění}\\
{\small vykonat, aplikovat, využít, demonstrovat v praxi, implementovat, vyřešit uzavřený a ohraničený problém}

\item \textbf{Analyzovat (\textit{analyze})}\\
\note{rozklad na části, vztahy a interakce mezi částmi}\\
{\small diskutovat, \enquote{rozbít} na části, rozlišit, navrhnout/zvolit řešení, organizovat, strukturovat, integrovat, propojit}

\item \textbf{Vyhodnotit (\textit{evaluate})}\\
\note{posouzení na základě stanovených kritérií a standardů}\\
{\small zkontrolovat, testovat, posoudit, odborně kritizovat, argumentovat}

\item \textbf{Vytvořit (\textit{create})}\\
\note{vytvoření nového celku, reorganizace do nové struktury}\\
{\small generovat, produkovat, vymyslet, plánovat/projektovat, komponovat, navrhovat, budovat, kreativně kombinovat prvky}
\end{enumerate}

\noteblack{Hranice mezi úrovněmi nejsou striktně určené. Seznam činností ti pomůže popsat cíle hodiny i znalosti a dovednosti, které by měli studenti získat. Ukáže ti také, jestli je tvá výuka dostatečně variabilní.}

\chapter*{Užitečné výukové nástroje}
\vspace{-0.5em}
\note{Seznam několika užitečných nástrojů a konceptů do výuky}

\section*{Výuka jako celek}

\begin{itemize}[leftmargin=1.5em]
\item Používání jmen -- \note{tykání, jmenovky, tahák se jmény}
\item Uspořádání prostoru -- \note{tvar usazení, viditelnost}
\item Verbalizace očekávání -- \note{očekávají obě strany to stejné?}
\item Precedenty -- \note{společná historie, co se může opakovat}
\item Zakázky -- \note{co studenti chtějí dostat?}
\item Hospitace -- \note{jak učí moji kolegové?}
\item Žádost o zpětnou vazbu -- \note{je to skutečně tak, jak si myslím?}
\item Bloomova taxonomie -- \note{verbalizace cílů hodiny/výuky}
\item Rubrika -- \note{sebehodnocení, sledování pokroku}
\end{itemize}

\section*{Struktura hodiny}

\begin{itemize}[leftmargin=1.5em]
\item Situování -- \note{shrnutí, kde jsme a kam jdeme}
\item Označení přechodů -- \note{explicitní přechod do dalšího bloku}
\item Veřejný checklist -- \note{viditelná struktura hodiny}
\item Check-in -- \note{otevření, zaměření pozornosti na tady a teď}
\item Check-out -- \note{uzavření, výzva k reflexi}
\item Tracking -- \note{sledování/uvědomování si vývoje diskuze/hodiny}
\item Parkoviště otázek -- \note{odkládání komplexních otázek na později}
\end{itemize}

\section*{Zadávání pokynů a úkolů}

\begin{itemize}[leftmargin=1.5em]
\item Otázky ke skupině -- \note{\enquote{kto} + podmínka + akce}
\item Multiple choice questions -- \note{uvěřitelnost možností}
\item Hlasování -- \note{pozitivní, negativní, rozdělování bodů, Kahoot}
\item Zadání úkolů -- \note{první krok, produkt/výsledek, hodnocení, trvání}
\item Praktická ukázka -- \note{zapojování vícero smyslů}
\item Rozcvička -- \note{fyzické/myšlenkové probuzení, opakování}
\item Nastavování prahu -- \note{jak moc snadné je zapojit se?}
\item Chunking -- \note{sdružování/dělení na malé \enquote{stravitelné} části}
\item Externí motivace -- \note{body navíc, odměny, bonbóny, \ldots}
\end{itemize}

\section*{Tvorba nových aktivit/cvičení}

\begin{itemize}[leftmargin=1.5em]
\item Schody -- \note{postup od lehčích úkolů k těžším}
\item Méně tradiční zadání -- \note{např.\ programování na mnoho způsobů}
\item Argumentace -- \note{zdůvodnění svého názoru, hodnocení jiných}
\item Figurky -- \note{fyzické objekty a manipulace s nimi}
\item Antiproblém -- \note{řešení opačného problému}
\item Myšlenkový oblak -- \note{sbírání asociací, diskuze, souvislosti}
\item Konceptová mapa -- \note{vizualizace pojmů/znalostí a jejich vztahů}
\end{itemize}

\section*{Uvádění aktivit}

\begin{itemize}[leftmargin=1.5em]
\item Skupinová práce -- \note{dělení skupiny na menší části/týmy}
\item Harvest/sklizeň -- \note{sdílení názorů, zesílení důležitých signálů}
\item Think-pair-share -- \note{nejprve sami, diskuze v páru, sdílení}
\item Peer review -- \note{studenti si navzájem hodnotí práci}
\item Peer tutoring -- \note{studenti se učí navzájem}
\item Neupozornění na chybu -- \note{studenti si ji najdou sami}
\item Obrácená hodina -- \note{samostudium teorie, společné cvičení}
\end{itemize}

\chapter*{Prostor pro poznámky}
\note{Např. nejdůležitější zpětná vazba z hospitací.}

\chapter*{Prostor pro poznámky}
\note{Např.: na co se chci soustředit, když učím?}

\chapter*{Prostor pro poznámky}
\note{Např.: ve kterých učitelských dovednostech jsem se zlepšil(a)?}

\chapter*{}
\putat{-22}{-372}{\includegraphics[width=\paperwidth]{../img/stack-ad-cs}}

% ===== IMPRINT =====

\newpage
\vspace*{\fill}
Reflektivní deník učitele

Verze: \today

\bigskip
Budeme velice rádi, když nám sdělíš svoje návrhy na zlepšení nebo postřehy z používaní deníku. Napiš nám email na adresu \textit{teachinglab@fi.muni.cz}.

Zdrojové soubory i PDF k tisku jsou dostupné online na adrese \url{github.com/teaching-lab/reflective-diary}.

\medskip
\begin{minipage}{0.28\textwidth}
	\includegraphics[width=0.80\textwidth]{../img/muni}\hfill
	\end{minipage}
\begin{minipage}{0.72\textwidth}
Deník vytvořili Martin Ukrop, Ondráš Přibyla, Valdemar Švábenský a Martin Macák pro kurz \textit{Teaching Lab} na Fakultě informatiky Masarykovy univerzity v spolupráci s Centrem jazykového vzdělávání Masarykovy univerzity.
\end{minipage}

\medskip
Naše zkušenosti z používání deníku jsou popsané v publikaci: M. Ukrop, V. Švábenský, J. Nehyba, \enquote{Reflective Diary for Professional Development of~Novice~Teachers}, in \textit{Proceedings of the 50th ACM Technical Symposium on Computer Science Education (SIGCSE)}, 2019.

\begin{minipage}{0.72\textwidth}
Toto dílo podléhá licenci Creative Commons Uveďte původ-Neužívejte dílo komerčně-Zachovejte licenci 4.0 Mezinárodní licence.
\end{minipage}
\begin{minipage}{0.28\textwidth}
\hfill\includegraphics[width=0.90\textwidth]{../img/cc-by-nc-sa}
\end{minipage}

% ===== BACK COVER PAGE =====
\newpage
\thispagestyle{empty}
\mbox{}
\newpage
\thispagestyle{empty}
\putat{-21}{-390}{\includegraphics[width=\paperwidth]{../img_cover/cover-back-cs}}

\end{document}
